\documentclass[12pt]{article}

\usepackage{sbc-template}

\usepackage{graphicx,url}
\usepackage{pgfplots}

\usepackage[brazil]{babel}
\usepackage[utf8]{inputenc}

\sloppy

\title{Identificação de Caminhos Disjuntos em Grafos Direcionados utilizando Fluxo Máximo}

\author{Adriano Araújo Domingos dos Santos\inst{1}}

\address{Instituto de Ciências Exatas e Informática -- Pontifícia Universidade
    \\ Católica de Minas Gerais (PUC Minas)
    \\ R. Cláudio Manoel, 1162 -- Savassi - 30.140-100 -- Belo Horizonte -- MG -- Brazil
    \email{\{adriano.domingos\}@sga.pucminas.br}
}

\begin{document} 

\maketitle

\section{Implementação}

A solução desenvolvida consiste na implementação do algoritmo de Dinic para cálculo de fluxo máximo em grafos direcionados, estendido para a identificação de caminhos disjuntos em arestas entre um par origem-destino. A aplicação foi desenvolvida em Java sem nenhuma dependência externa, salvo para os testes unitários que não são necessários para sua execução.

\subsection{Algoritmo de Dinic} \label{dinic}

O algoritmo de Dinic foi implementado em sua formulação clássica, composta por duas fases principais: a construção de um grafo de níveis (\textit{Level Graph}) por meio de uma BFS a partir da origem, e a busca de caminhos aumentantes no grafo de níveis através de uma DFS (\textit{Blocking Flow}). A estrutura residual do fluxo é armazenada em vetores planos de inteiros para capacidade residual (\texttt{cap}), fluxo atual (\texttt{flo}) e vértice destino (\texttt{to}). Cada aresta original $i$ é mapeada para dois índices consecutivos na estrutura residual: $2i$ para a aresta direta e $2i+1$ para a aresta reversa, permitindo a atualização do fluxo em ambos os sentidos utilizando o operador XOR. A otimização de ponteiro atual (vetor \texttt{it}) é empregada na DFS para evitar a varredura repetida de arestas já saturadas, garantindo a complexidade $\mathcal{O}(E \sqrt{V})$ para o caso de capacidades unitárias.

\subsection{Caminhos Disjuntos} \label{disjointpaths}

A identificação de caminhos disjuntos em arestas é realizada através de uma redução para o problema de fluxo máximo. Dado um grafo direcionado $G = (V, E)$ e um par origem-destino $(s, t)$, cada aresta recebe capacidade unitária, transformando $G$ em uma rede de fluxo. O algoritmo de Dinic é então executado para encontrar o fluxo máximo de $s$ para $t$. Como cada aresta possui capacidade $1$, o valor do fluxo máximo encontrado corresponde exatamente ao número máximo de caminhos disjuntos em arestas entre $s$ e $t$. Após o cálculo do fluxo, os caminhos são extraídos percorrendo-se o grafo a partir de $s$, seguindo arestas com fluxo positivo, e marcando arestas já utilizadas para garantir a disjunção entre as rotas. Ao final, a implementação exibe a quantidade de caminhos disjuntos encontrados e lista individualmente cada caminho como uma sequência de vértices da origem ao destino.

\subsection{Representação \textit{Forward-Star}} \label{forwardstar}

A representação \textit{Forward Star} é uma estrutura de dados baseada em vetores contíguos, que serializa a topologia do grafo em três vetores primitivos unidimensionais: \texttt{pointers}, \texttt{targets} e \texttt{weights}. Ela foi escolhida por conta de seu acesso rápido à adjacência dos vértices e excelente uso da localidade de cache.

Na implementação desenvolvida, a construção do grafo foi estritamente desacoplada de sua representação final através do padrão de projeto \textit{Builder}. Durante a fase de construção, as arestas são inseridas em vetores temporários paralelos. Ao finalizar a instanciação, o construtor ordena as arestas por vértice de origem, compacta os destinos no vetores finais \texttt{targets} e \texttt{weights} de tamanho $|E|$, e preenche o vetor \texttt{pointers} de tamanho $|V| + 1$. Cada índice $v$ em \texttt{pointers} armazena o índice inicial de seus vizinhos no vetor \texttt{targets}, permitindo que a varredura da vizinhança $\Gamma(v)$ ou dos sucessores $\Gamma^+(v)$ de qualquer vértice ou o cálculo do seu grau sejam realizados em tempo $\mathcal{O}(1)$.

\subsection{Gerador de Grafos Sintéticos} \label{gerador}

Para viabilizar a análise de desempenho em cenários massivos e controlados, foi desenvolvida uma ferramenta própria de geração de grafos sintéticos. O gerador é capaz de instanciar topologias direcionadas e não direcionadas, parametrizadas pela quantidade exata de vértices ($|V|$), arestas ($|E|$) e requisitos de conectividade.

A fim de satisfazer a conectividade exigida (Simplesmente ou Fortemente Conexo), a estratégia de construção foi dividida em três etapas fundamentais:

\begin{enumerate}
    \item \textbf{Garantia de Conectividade Base:} Inicialmente, um vetor contendo todos os identificadores de vértices é embaralhado de forma aleatória utilizando o algoritmo de \textit{Fisher-Yates}. Em seguida, itera-se sobre esse vetor conectando o vértice $i$ ao $i+1$, formando um caminho Hamiltoniano que garante a alcançabilidade de todos os vértices. Para grafos fortemente conexos, uma aresta de fechamento (do último vértice para o primeiro) é adicionada, transformando o caminho em um ciclo.
    
    \item \textbf{Preenchimento de Densidade:} Para atingir a cardinalidade $|E|$ desejada, o algoritmo entra em um laço de amostragem uniforme, sorteando pares de vértices $(u, v)$ aleatoriamente. Para assegurar que o grafo resultante seja um grafo simples (sem arestas múltiplas), as conexões geradas são validadas contra um conjunto de dispersão (\texttt{HashSet}), que armazena as arestas para buscas em tempo $\mathcal{O}(1)$ utilizando empacotamento de bits para diminuir o custo de armazenamento.
    
    \item \textbf{Ponderação:} Após a definição da topologia, a matriz de pesos é populada. Cada aresta recebe um valor inteiro sorteado mediante uma distribuição uniforme dentro do limite parametrizado pelo usuário.
\end{enumerate}

Essa abordagem garante que as instâncias geradas para o experimento sejam conexas, livres de laços e arestas paralelas.

\section{Metodologia} \label{metodologia}

Para avaliar a eficácia e a eficiência da implementação, foi desenvolvido um script automatizado de \textit{benchmarking}. O experimento foi desenhado para testar os limites da estrutura e do algoritmo utilizando concorrência (\textit{threads}) para acelerar a execução das múltiplas tentativas.

\subsection{Design do Experimento e Parâmetros}

Para garantir a significância estatística e mitigar anomalias causadas por processos de \textit{background} do sistema operacional ou pelo \textit{Garbage Collector} da JVM, cada cenário único foi executado em 10 repetições independentes. Os parâmetros estabelecidos foram:

\begin{itemize}
    \item \textbf{Topologia do Grafo:} Foram avaliados grafos sintéticos direcionados sob duas conectividades: \textit{Fortemente Conexo} e \textit{Euleriano}.
    
    \item \textbf{Tamanho do Grafo:} Foram gerados grafos em quatro escalas crescentes de magnitude. Os pares de vértices ($|V|$) e arestas ($|E|$) utilizados foram: $(1.000, 100.000)$, $(20.000, 2.000.000)$, $(500.000, 50.000.000)$ e $(1.000.000, 100.000.000)$.
\end{itemize}

A combinação ortogonal de todas essas variáveis (2 tipos $\times$ 4 tamanhos) executada 10 vezes resultou em um espaço amostral de 80 execuções documentadas.

\subsection{Fluxo de Execução e Coleta de Dados}

Para cada iteração de teste, o fluxo seguiu um controle rigoroso para isolar os tempos:

\begin{enumerate}
    \item \textbf{Geração:} Instanciação do grafo sintético.
    
    \item \textbf{Construção da Estrutura:} Transcrição das arestas geradas para a estrutura de dados utilizando a representação \textit{Forward Star}.
    
    \item \textbf{Transformação do Grafo:} Atribuição da capacidade de cada aresta igual a 1.

    \item \textbf{Busca dos Distâncias:} Execução isolada do algoritmo de Dinic, encontrando o fluxo máximo do vértice 1 até o vértice $n$.
    
    \item \textbf{Registro:} Coleta e exportação das métricas de desempenho.
\end{enumerate}

\subsection{Métricas de Avaliação}

As métricas coletadas incluíram o tempo de geração do grafo, o tempo de processamento efetivo do algoritmo de Dinic, quantidade do caminhos disjuntos encontrados e o tipo do grafo.

\section{Resultados} \label{resultados}

Os experimentos foram conduzidos conforme a metodologia descrita na Seção \ref{metodologia}. A Tabela \ref{tab:resultados} apresenta os resultados consolidados para os oito cenários avaliados, com médias calculadas sobre 10 execuções independentes.

\begin{table}[ht]
\centering
\caption{Desempenho médio do Dinic por tipo de conectividade.}
\label{tab:resultados}
\resizebox{\textwidth}{!}{
\begin{tabular}{cccccc}
\hline
$|\textbf{V}|$ & $|\textbf{E}|$ & \textbf{Conectividade} & \textbf{Caminhos} & \textbf{Tempo (ms)} \\
\hline
1.000     & 100.000     & Fortemente Conexo & 92,90 & 62,30 \\
20.000    & 2.000.000   & Fortemente Conexo & 93,80 & 2251,00 \\
500.000   & 50.000.000  & Fortemente Conexo & 95,70 & 70.550,30 \\
1.000.000 & 100.000.000 & Fortemente Conexo & 97,60 & 155.477,90 \\
\hline
1.000     & 100.000     & Euleriano         & 93,60 & 45,40 \\
20.000    & 2.000.000   & Euleriano         & 95,00 & 2334,30 \\
500.000   & 50.000.000  & Euleriano         & 98,10 & 76.819,90 \\
1.000.000 & 100.000.000 & Euleriano         & 91,10 & 154.061,90 \\
\hline
\end{tabular}
}
\end{table}

Os resultados consolidados na Tabela \ref{tab:resultados} demonstram a escalabilidade da implementação do algoritmo de Dinic para grafos de grande porte. O tempo de execução apresentou crescimento aproximadamente linear em escala log-log (Figura \ref{fig:line-growth}), consistente com a complexidade esperada $\mathcal{O}(E \sqrt{V})$ para capacidades unitárias. Para a instância de maior magnitude (1 milhão de vértices e 100 milhões de arestas), o tempo médio de execução foi de aproximadamente 155 segundos para ambas as conectividades, evidenciando a viabilidade prática da abordagem mesmo em cenários massivos.

\begin{figure}[htbp]
    \centering
    \begin{tikzpicture}
        \begin{axis}[
            xmode=log,
            ymode=log,
            xlabel={Número de Vértices ($|V|$) -- Escala Logarítmica},
            ylabel={Tempo Médio (ms) -- Escala Logarítmica},
            legend pos=north west,
            grid=major,
            grid style={dashed, gray!30},
            width=\textwidth,
            height=8cm,
            log basis x={10},
            log basis y={10},
            mark size=2.5pt
        ]
        
        \addplot[color=blue, mark=square*, thick] coordinates {
            (1000, 62.30 )
            (20000, 2251.00 )
            (500000, 70550.30 )
            (1000000, 155477.90 )
        };
        \addlegendentry{Fortemente Conexo}
        
        \addplot[color=red, mark=triangle*, thick] coordinates {
            (1000, 45.40)
            (20000, 2334.30)
            (500000, 76819.90)
            (1000000, 154061.90)
        };
        \addlegendentry{Euleriano}
        
        \end{axis}
    \end{tikzpicture}
    \caption{Crescimento logarítmico do tempo de execução do Dinic em relação ao tamanho da entrada.}
    \label{fig:line-growth}
\end{figure}

Em relação à quantidade de caminhos disjuntos encontrados (coluna Caminhos da Tabela \ref{tab:resultados}), observa-se que ambos os tipos de conectividade apresentaram valores próximos ao grau médio do grafo ($|E|/|V| \approx 100$), indicando alta capilaridade da topologia gerada pelo algoritmo de Fisher-Yates com preenchimento aleatório. O número de caminhos manteve-se estável ao longo das diferentes escalas, sugerindo que a estrutura dos grafos sintéticos proporciona diversidade de rotas alternativa entre origem e destino independentemente do tamanho da instância.

\section{Conclusão}

A implementação do algoritmo de Dinic combinada com a representação \textit{Forward Star} demonstrou capacidade de processar grafos com até 100 milhões de arestas em tempos viáveis para análise experimental. A redução do problema de caminhos disjuntos em arestas para fluxo máximo com capacidades unitárias mostrou-se viável, recuperando todos os caminhos sem violar a restrição de disjunção. Os resultados confirmam a escalabilidade da abordagem, habilitando sua aplicação para a identificação dos caminhos disjuntos.

\end{document}
